% sections/01-introduction.tex
\section{Introduction}\label{sec:intro}

We prove that spectral-gap certificates for the Laughlin $V_1$ model pass from the plane to rectangular tori with explicit
errors. The main result is a transfer principle (Theorem~\ref{thm:transferintro}): every planar certificate of the form introduced
by OpenAI~\cite{OAI:A-Fock-space-inequality-and-the-Laughlin-spectral-gap-September-24-2026}, with constant $\gamma$, gives
$\HT^2\ge(\gamma-\etr(\Tor))\HT$ on the whole Fock space for the untruncated $V_1$ Hamiltonian $\HT$ of every rectangular
torus, where $\etr(\Tor)$ is bounded by an explicit function of the shorter period that tends to zero at a Gaussian rate. For
the certificate of~\cite{OAI:A-Fock-space-inequality-and-the-Laughlin-spectral-gap-September-24-2026} we evaluate these bounds.
On every rectangular torus whose shorter period is at least $L$, with $L\ge16$, the inequality $\HT^2\ge\gT\HT$ holds with a
constant $\gT\ge\gl(L)$, uniformly in the particle number and in the aspect ratio, where $\gl$ is an explicit nondecreasing
function that tends to the planar constant $\gs$
of~\cite{OAI:A-Fock-space-inequality-and-the-Laughlin-spectral-gap-September-24-2026} at a Gaussian rate in $L$; both
$\gs-\gl(L)$ and the bound on $\etr(\Tor)$ are $O(L^4\,2^{-L^2/20})$ (Remark~\ref{rem:rate}).
The thresholds are explicit: a numerical evaluation of $\gl$ (Remark~\ref{rem:evaluation}) gives $\gl(L)>0$ for $L\ge16.457$,
$\gl(L)>1/25$ for $L\ge16.857$ and $\gl(L)\ge\gs-10^{-5}$ for $L\ge18.685$, in units of the magnetic length
(Theorem~\ref{thm:openai}).

Laughlin's wave function~\cite{Laughlin} is the exact zero-energy ground state of Haldane's $V_1$ pseudopotential
Hamiltonian~\cite{Haldane}, in which a pair of spinless fermions in the lowest Landau level pays energy one if its relative
angular momentum is one and nothing otherwise. Whether this Hamiltonian has a spectral gap above its zero modes that stays
bounded below as the system grows is a long-standing question~\cite[Appendix~A]{Rougerie}. For the untruncated interaction it
was answered on the sphere and in the plane by OpenAI~\cite{OAI:A-Fock-space-inequality-and-the-Laughlin-spectral-gap-September-24-2026}:
with $\gs=4616733319001/10^{14}>1/25$, the spherical Hamiltonian $H_Q$ satisfies $H_Q^2\ge\gamma H_Q$ on the whole Fock space for every
$\gamma<\gs$ and all sufficiently large flux $Q$, and the planar Hamiltonian satisfies $(H^{\mathrm{pl}}_N)^2\ge\gs H^{\mathrm{pl}}_N$ for every
particle number. The proof compares $H_Q^2$ with a positive operator built from seven explicit rows, averaged over rotations; the
rows satisfy finitely many planar inequalities, verified by exact rational arithmetic. Uniform gaps for truncated versions of the
pseudopotential were obtained earlier, on thin cylinders and thin tori~\cite{NWY,WY}. Two further works treat the untruncated
interaction on the sphere: Gao~\cite{Gao} raised the constant to about $0.219$, on the sphere and in the plane, by a semidefinite
search, and Scheer and Khalaf~\cite{ScheerKhalaf} constructed a semidefinite hierarchy of lower bounds on the sphere, with exact
rational certificates and a rigorous lower bound on the limit inferior of the gap as the flux grows, and leave the torus to
future work. In all three spherical results the flux threshold is existential: it is shown to exist, not computed.

A uniform bound of this kind, for every flat torus whose shortest lattice vector exceeds an existential threshold, has also
been claimed in a recent candidate proof prepared by Hou with advice from Wu~\cite{HouWu}; the present argument, obtained
independently, is restricted to rectangular tori but makes the threshold, the constant $\gl(L)$ and its Gaussian rate of
convergence to $\gs$ explicit. That candidate proof, which is unrefereed and not on arXiv, states the bound $H^2\ge\gamma H$ for
every $\gamma<\gs$, uniformly in the particle number, on tori of every lattice shape and with every boundary twist, and also
claims that the kernel at filling one third is three-dimensional for every shape. Where the two overlap --- the uniform
bound on rectangular tori and the three-dimensional kernel --- the statements agree. It transfers fixed certificates and states
no principle for a class of certificates; Theorem~\ref{thm:transferintro} below is such a principle, with its error explicit.

\subsection{The transfer principle}\label{ssec:transfer}

Let $\Tor$ be a rectangular torus with periods $L_x,L_y$ (in units of the magnetic length) and $\Nphi=L_xL_y/2\pi$ flux quanta,
let $\h$ be its lowest Landau level and $\Fock$ the fermionic Fock space over $\h$, and let $\HT$ be the untruncated $V_1$
Hamiltonian: the second quantization of the periodization, over the lattice of periods, of the planar projector onto pairs of
relative angular momentum one (Definition~\ref{def:HT}). Let $L:=\min\{L_x,L_y\}$ be the shorter period.

The certificate of~\cite{OAI:A-Fock-space-inequality-and-the-Laughlin-spectral-gap-September-24-2026} is a finite family of rows
together with planar inequalities that they satisfy. On the torus, at a base point $R$, row $\nu$ is the operator
\[
F_\nu(R)=\lambda_\nu\,B(\ut_{t_\nu}(R))+\sum_{(p,j)\in\mathcal S_\nu}\alpha^\nu_{pj}\,c(\chi^R_{p+j-t_\nu})^*\,c(\chi_j^R)\,B(\ut_p(R))
\]
of Definition~\ref{def:KT}, where $B(\ut_p(R))$ removes a pair in the periodized pair state of label $p$ about $R$ and $c(\chi^R_m)$ a
particle in the periodized orbital of degree $m$ about $R$ (Section~\ref{sec:local}). The proof uses the certificate only through
the data of the rows and the inequalities, which define the following class.

\begin{definition*}[planar certificates of the same form]
A \emph{planar certificate of the same form} as that
of~\cite{OAI:A-Fock-space-inequality-and-the-Laughlin-spectral-gap-September-24-2026} consists of finitely many rows $\nu$, each with an
integer $t_\nu\ge0$, a real $\lambda_\nu$ and real coefficients $\alpha^\nu_{pj}$ supported in a finite set
$\mathcal S_\nu\subset\{(p,j):0\le p\le P_c,\ 0\le j\le M_c,\ 0\le p+j-t_\nu\le M_c\}$, where $P_c,M_c\ge0$ are integers, and an
allowance $\eps\ge0$, such that, with the cutoffs $Z_c:=\max\{p+j:(p,j)\in\bigcup_\nu\mathcal S_\nu\}$ and
$D_c:=\max\{p+j+(q+l-t_\nu)\}$ over $(p,j),(q,l)\in\mathcal S_\nu$, the inequalities
\begin{equation}\label{eq:certgeneral}
e_z\le\tfrac32q_z\ \ (z\le Z_c,\ z\notin\{0,1,3\}),\qquad G_D(2\kb{e_1}{e_1}-E_D+\eps\one)G_D\ge0\ \ (1\le D\le D_c)
\end{equation}
hold. Here $q_z=(3z-1)(-\frac12)^z$ are the planar three-body coefficients (Lemma~\ref{lem:qz}); $e_z$ and $E_D$ are the three-body
scalars and the four-body copy matrices of the rows, defined by \eqref{eq:ez} and \eqref{eq:ED} from the planar coupling
coefficients (the sum over $T$ in \eqref{eq:ez} running to $Z_c$, and \eqref{eq:ED} taken for $1\le D\le D_c$); $G_D$ is the Gram
matrix of the planar four-body relative states $W_4\omega_{D,r}$, $r\in\mathcal R_D$, defined with them in Section~\ref{sec:moyal}; $\mathcal R_D$
is the set of odd integers in $[1,D]$; and $\kb{e_1}{e_1}$ is the matrix with $(1,1)$ entry $1$ and all other entries zero. The
\emph{constant} of the certificate is
\[
\gamma:=1-\etas-\delta_\infty(Z_c)-\eps\,N_4,\qquad \etas=\sum_\nu\lambda_\nu^2,
\]
\[
\delta_\infty(Z_c):=\frac32\sum_{\substack{z>Z_c\\ z\ \mathrm{odd}}}\frac{(3z-1)(z+1)}{2^z},\qquad N_4:=\sum_{D=1}^{D_c}\Bigl\lfloor\frac{(D+1)^2}4\Bigr\rfloor .
\]
We require that some $\mathcal S_\nu$ be nonempty, that $M_c\ge1$ and that $D_c\ge2$, so that $Z_c$ is defined and the bounds of
Lemma~\ref{lem:Ebounds}(b) are used only for $m\ge1$; the certificate of~\cite{OAI:A-Fock-space-inequality-and-the-Laughlin-spectral-gap-September-24-2026}
has $M_c=8$ and $D_c=23$.
\end{definition*}

Every row of this form ends in a pair annihilator and has a definite weight $t_\nu+1$ (each of its summands removes two particles
and lowers the total orbital degree by $t_\nu+1$), and these two features are all that the proof uses of a row; that is why the
principle is a property of the form, not of the numbers in the certificate
of~\cite{OAI:A-Fock-space-inequality-and-the-Laughlin-spectral-gap-September-24-2026}.

\begin{theorem}[the transfer principle]\label{thm:transferintro}
Given a planar certificate of the same form with constant $\gamma$, for every rectangular torus, on $\Fock$,
\[
\HT^2\ge(\gamma-\etr(\Tor))\,\HT,
\]
where, with $\Ec_m=\Ec_m(\Tor)$ and $\sigma_z=\sigma_z(\Tor)$,
\begin{multline*}
\etr(\Tor):=(1-\lambda_{\min})+\frac32\sum_{\substack{z>Z_c\\ z\ \mathrm{odd}}}\frac{(3z-1)(z+1)}{2^z}\sigma_z
+(M_c+1)(P_c+1)\sum_\nu\norm{\alpha^\nu}^2\,\Ec_{M_c}(1+\Ec_{M_c})\\
+\eps N_4\bigl((1+\Ec_{D_c})^2-1\bigr)+2N_4\,\Ec_{D_c-1}(1+\Ec_{D_c})^2 .
\end{multline*}
Here $\lambda_{\min}$ is the smallest eigenvalue of the Gram matrix of the torus pair functions (Lemma~\ref{lem:levels}),
$\sigma_z(\Tor)=\sup_R\max_{m\le z}\bigl(\norm{\chi^R_m}^2-1\bigr)$ (Lemma~\ref{lem:Ebounds}(c)), $\Ec_m(\Tor)$ is the lattice sum, over
the nonzero periods, of the norms of the truncated displacement matrices $\Dis^{(m)}$ (Section~\ref{sec:local}), and
$\norm{\alpha^\nu}^2=\sum_{p,j}(\alpha^\nu_{pj})^2$. For every rectangular torus whose shorter period is at least $L$, $\etr(\Tor)\le\etr(L)$
with $\etr(L)\to0$ as $L\to\infty$. In particular, if $\gamma>0$, then for every $\gamma'<\gamma$ there is $L_0$ such that $\HT^2\ge\gamma'\HT$
on the whole Fock space of every rectangular torus whose shorter period is at least $L_0$, uniformly in the particle number and the
aspect ratio.
\end{theorem}

The function $\etr(L)$, obtained in the proof from explicit bounds on the terms of $\etr(\Tor)$, is $O(L^4\,2^{-L^2/20})$ as
$L\to\infty$, with a constant that depends on the certificate (Remark~\ref{rem:rate}). The theorem does not require the planar
certificate to come from a spherical proof: any finite family of rows of this form satisfying \eqref{eq:certgeneral} with $\gamma>0$
yields a gap on rectangular tori, uniform above a threshold period, with that rate.

For the certificate of~\cite{OAI:A-Fock-space-inequality-and-the-Laughlin-spectral-gap-September-24-2026}
(Section~\ref{sec:moyal} and Appendix~\ref{app:certificate}), $P_c=7$, $M_c=8$, $Z_c=15$, $D_c=23$, $\delta_\infty(15)=61/4096$ and
$N_4=1222$, with $\etas=93527408868499/10^{14}$ and $\eps=3/10^6$; its constant is $\gamma=\gs=4616733319001/10^{14}>1/25$, and
$\gs-\etr(\Tor)$ is, term by term, the constant $\gT$ of \eqref{eq:gT}. Theorem~\ref{thm:transferintro} for this certificate, together
with an evaluation of $\etr$ for it, gives the following explicit bounds.

\begin{theorem}[OpenAI's certificate on rectangular tori]\label{thm:openai}
For every rectangular torus $\Tor$, on the whole Fock space,
\[
\HT^2\ge\gT\,\HT,
\]
with the explicit constant $\gT$ of \eqref{eq:gT}, which depends on the torus only through its lattice of periods. For $L\ge16$ let
$\gl(L)$ be the explicit function of Corollary~\ref{cor:uniform}. Then $\gT\ge\gl(L)$ for every rectangular torus whose shorter period
is at least $L$, so that $\HT^2\ge\gl(L)\HT$ uniformly in the particle number and in the aspect ratio; $\gl$ is nondecreasing on
$[16,\infty)$, and $\gl(L)\to\gs$ as $L\to\infty$, with $\gs-\gl(L)=O(L^4\,2^{-L^2/20})$. Numerically, $\gl(L)>0$ for $L\ge16.457$,
$\gl(L)>1/25$ for $L\ge16.857$ and $\gl(L)\ge\gs-10^{-5}$ for $L\ge18.685$, in units of the magnetic length; these values are
floating-point evaluations of rigorous closed-form bounds, not interval-arithmetic certificates (Remark~\ref{rem:evaluation}).
\end{theorem}

Theorem~\ref{thm:openai} is thus the case of Theorem~\ref{thm:transferintro} for the certificate
of~\cite{OAI:A-Fock-space-inequality-and-the-Laughlin-spectral-gap-September-24-2026}, with $\etr$ evaluated. Section~\ref{sec:theorem}
proves the two in the opposite order: Theorem~\ref{thm:main} and Corollary~\ref{cor:uniform}, with Remark~\ref{rem:rate} and the
thresholds evaluated in Section~\ref{ssec:constants}, prove Theorem~\ref{thm:openai} directly, and Section~\ref{ssec:transferproof} then
proves Theorem~\ref{thm:transferintro}, restated as Theorem~\ref{thm:transfer}, by the same argument with the data of the certificate
left general.

\begin{remark}[scope]\label{rem:scope}
Theorem~\ref{thm:transferintro} is proved for the form of the definition: rows each of whose summands is a pair annihilator, alone
or dressed by one creation and one annihilation operator; the two-body $V_1$ interaction of a single species of spinless fermions;
and rectangular tori without boundary twists. Three natural extensions are not covered. Rows of higher degree in the fermion
operators have squares with parts on five or more particles; they would need the torus Moyal identity on that many particles,
which Theorem~\ref{thm:moyal} states and proves for every number of particles, and a comparison term for those parts, which is not
constructed here. Two species, as in the spinful hollow-core model of Section~\ref{sec:outlook}, and general lattices of periods
with boundary twists are not treated. Whether these extensions hold is not claimed here.
\end{remark}

By Theorem~\ref{thm:openai}, the spectrum of $\HT$ in every particle sector lies in $\{0\}\cup[\gT,\infty)$. At $\Nphi=3N$ the kernel on
$N$ particles is three-dimensional, spanned by the three centre-of-mass copies of the Laughlin state, and no sector with more than
$N$ particles has a zero mode (Theorem~\ref{thm:kernel}, Proposition~\ref{prop:overfill}); the same count holds, for every modular
parameter, for the explicit theta-function pair operators $B_n^{(\tau)}$ of \eqref{eq:Bntau}, whose identification with the $V_1$
Hamiltonian of an oblique torus is not proved here. So $\HT\ge\gT(\one-P_N)$ on $N$ particles, with $P_N$ of rank three, and
$\HT\ge\gT$ on more particles (Corollary~\ref{cor:spectrum}).

At the square aspect ratio $L=\sqrt{2\pi\Nphi}$, so the bound is positive for $\Nphi\ge44$ and exceeds $1/25$ for $\Nphi\ge46$. At
filling one third ($\Nphi=3N$) this means a gap for $N\ge15$ ($\gl(L)\ge0.0385$ at $N=15$) and a gap larger than $1/25$ for $N\ge16$
($\gl(L)\ge0.0456$ at $N=16$). For elongated tori the condition is on the shorter period, so larger systems are needed. All error terms
decay like $e^{-cL^2}$ in the shorter period; the binding one at the thresholds is a lattice sum of norms of truncated
displacement matrices, bounded by a multiple of $L^{44}e^{-L^2/4}$ (Lemma~\ref{lem:Ebounds}), and the thresholds are where it
becomes small. The slowest-decaying bound is that on the three-body tail beyond the blocks covered by the certificate,
$O(L^4\,2^{-L^2/20})$ (Proposition~\ref{prop:three}); it sets the rate of Remark~\ref{rem:rate}.

\subsection{From the sphere to the torus}

The torus proof follows the architecture
of~\cite{OAI:A-Fock-space-inequality-and-the-Laughlin-spectral-gap-September-24-2026}. One expands $\HT^2$ in normal order into
two-, three- and four-body terms, and compares them with the corresponding terms of a positive comparison operator
$\KT=\frac1\pi\int_\Tor\sum_\nu F_\nu(R)^*F_\nu(R)\,\dd R$, where the rows $F_\nu(R)$ of Section~\ref{ssec:transfer} are the planar
rows of the certificate, those of~\cite{OAI:A-Fock-space-inequality-and-the-Laughlin-spectral-gap-September-24-2026} for
Theorem~\ref{thm:openai}, built at the base point $R$ from periodized orbitals and pair states. The planar certificate is used as
data, without change: for that of~\cite{OAI:A-Fock-space-inequality-and-the-Laughlin-spectral-gap-September-24-2026}, the row
coefficients, the thirteen three-body inequalities and the twenty-three four-body matrix inequalities. Section, lemma and equation numbers
of~\cite{OAI:A-Fock-space-inequality-and-the-Laughlin-spectral-gap-September-24-2026} refer to its version of 24 September 2026,
and its lemmas are also cited by name. On the sphere the comparison operator is
averaged over rotations and the finite-flux corrections are controlled by a relative-error transfer lemma; on the torus the
average over the base point $R$ replaces the rotations, and the corrections take a different and simpler form. Four facts,
proved here, make the transcription work:
\begin{enumerate}[(1)]
\item \emph{The projector replacement.} On two particles $\HT$ is not a projection. The Gram matrix of the torus pair functions is
computed in closed form as a lattice sum of image overlaps (Lemma~\ref{lem:gram}); its eigenvalues are
$1\pm\Dlev(L_y)\pm\Dlev(L_x)$ for even $\Nphi$ and $1\pm(\Dlev(L_x)^2+\Dlev(L_y)^2)^{1/2}$ for odd $\Nphi$, up to a smaller remainder, with
$\Dlev(\ell)=2(\ell^2/4-1)e^{-\ell^2/8}$ (Lemma~\ref{lem:levels}). Only the smallest eigenvalue enters the budget.
\item \emph{Periodized local modes.} The periodized orbitals about a point have an explicit Gram matrix, a lattice sum of
displacement matrices with phases (Lemma~\ref{lem:localgram}), whose deviation from the identity is bounded uniformly by lattice
sums of norms of truncated displacement matrices (Lemma~\ref{lem:Ebounds}).
\item \emph{Torus Moyal identities.} Averaging a projector onto periodized $k$-particle states with a fixed relative part and a
centre-of-mass orbital about $R$ over $R\in\Tor$ gives exactly the periodized image of the relative projector, and different
centre-of-mass degrees average to zero (Theorem~\ref{thm:moyal}). This replaces Schur's lemma for $SU(2)$; the proof uses Schur's
lemma for the magnetic translations of the centre-of-mass torus. With it, the pair-level part of $\KT$ is exactly $\etas\HT$, and
its three- and four-body parts are exactly the periodized images of the planar partial traces that define the certificate,
up to one small non-orthonormality term (Proposition~\ref{prop:parts}).
\item \emph{The square of $\HT$.} The three-body coefficient of $\HT^2$ is a sum of products of two pair operators sharing exactly
one particle, and periodization is multiplicative on such products (Lemma~\ref{lem:product}). Hence it is \emph{exactly} the
periodized planar three-body target, with the planar coefficients $q_z=(3z-1)(-\frac12)^z$ on every block
(Theorem~\ref{thm:square}). The three-body target therefore transfers without error, and the tail beyond the blocks covered by
the certificate is bounded with an explicit rate (Proposition~\ref{prop:three}). The four-body coefficient involves products on
disjoint pairs, and a Bessel-type bound for the localized channels of the relative coordinate of the two pair centres closes
the budget (Lemma~\ref{lem:bessel}).
\end{enumerate}
The resulting error terms are all relative to $\HT$ (they vanish on its kernel and carry no factor of the particle number), and all
are Gaussian in the shorter period. They are collected in the error $\etr(\Tor)$ of Theorem~\ref{thm:transferintro}, and for the
certificate of~\cite{OAI:A-Fock-space-inequality-and-the-Laughlin-spectral-gap-September-24-2026} in the constant $\gT$
of \eqref{eq:gT}.

Section~\ref{sec:kernel} determines the kernel at $\Nphi=3N$, for the pair operators $B_n^{(\tau)}$ of \eqref{eq:Bntau} at every
modular parameter and hence for $\HT$ on every rectangular torus. The threefold degeneracy of the torus Laughlin states is
classical~\cite{HaldaneRezayiTorus}; that these states are the only zero modes of the torus $V_1$ Hamiltonian is stated for
rectangular tori in~\cite[Section~II.2]{ONDS}, and a proof for every lattice shape and boundary twist is claimed
in~\cite{HouWu}. The argument of Section~\ref{sec:kernel} uses third-order vanishing at coincidences, division by the cube of
the theta-function Jastrow factor, and Liouville's theorem.

\subsection{Structure of the paper}

Section~\ref{sec:setting} sets up the torus, the planar lowest Landau level, the periodization and the periodized image of planar
operators, and defines $\HT$; Proposition~\ref{prop:orbital} gives its orbital form. Section~\ref{sec:kernel} determines the
kernel. Section~\ref{sec:projector} computes the pair Gram matrix and records the normal-ordered square of a lifted pair
operator. Section~\ref{sec:local} introduces the periodized local modes and bounds their Gram matrix. Section~\ref{sec:moyal}
proves the torus Moyal identities, states the certificate, and decomposes the comparison operator. Section~\ref{sec:square}
proves the exact three-body identity and the three-body comparison, and Section~\ref{sec:bessel} the four-body Bessel bound.
Section~\ref{sec:theorem} proves Theorem~\ref{thm:transferintro} and evaluates it for the certificate
of~\cite{OAI:A-Fock-space-inequality-and-the-Laughlin-spectral-gap-September-24-2026}: it assembles the budget and evaluates the
constants for that certificate, which proves Theorem~\ref{thm:openai}, derives the spectral consequences, and then proves
Theorem~\ref{thm:transferintro} by the same argument with the data of the certificate left general.
Section~\ref{sec:outlook} discusses the transfer principle as a basis for further work and lists what is not proved.
Appendix~\ref{app:certificate} records the certificate data, and Appendix~\ref{app:notation} lists the notation.
